\begin{textsample}{NEG Dim 2 – 2017\_c4 – Score 10.00 – 2017\_c4/enade\_2017\_eng\_mec\_6ddc12ff1b.txt}  \label{ex:f2_neg_342}
CAPÍTULO 4 PERCEPÇÃO DA PROVA As análises feitas neste capítulo tratam das percepções dos concluintes da Área de Engenharia Mecânica sobre a prova aplicada no Enade/2017.
Essas percepções foram mensuradas por meio de nove questões que avaliaram desde o grau de dificuldade da prova até o tempo gasto para concluí -la.
As percepções sobre a prova foram relacionadas ao desempenho dos estudantes e à Grande Região de funcionamento do curso.
O questionário de percepção da prova encontra -se ao final do Anexo VII, que traz a reprodução do exame.
O desempenho dos estudantes foi classificado em quatro quartos.
Para tanto, esse desempenho foi ordenado de forma ascendente.
O percentil 25, P25, também conhecido como primeiro quartil, é a nota de desempenho que deixa um quarto ( 25 \% ) dos valores observados abaixo e três quartos acima.
A Figura 1 apresenta uma ilustração deste conceito.
O quarto inferior de desempenho é composto pelas notas abaixo do primeiro quartil.
Já o percentil 75, P75, também conhecido como terceiro quartil, é o valor para o qual há três quartos ( 75 \% ) dos dados abaixo e um quarto acima dele.
O quarto superior de desempenho é composto pelas notas iguais ou acima do terceiro quartil.
O percentil 50, P50, também conhecido como mediana, é o valor que divide as notas em dois conjuntos de igual tamanho.
O segundo quarto inclui valores entre o primeiro quartil ( P25 ) e a mediana.
O terceiro quarto contém os valores entre a mediana ( P50 ) e o terceiro quartil ( P75 ).
Vale ressaltar que percentis, quartis e medianas são pontos que não obrigatoriamente pertencem ao conjunto original de dados, ao passo que os quartos são subconjuntos dos dados originais.
A seguir, serão apresentados gráficos com resultados selecionados, relativos às nove questões avaliadas por grupos de estudantes.
Os gráficos apresentam nas barras o percentual de alunos que assinalaram uma das opções ou a soma das porcentagens daqueles que assinalaram duas ( ou três ) delas.
Por exemplo, para as questões 1 e 2, os gráficos apresentam a porcentagem total de participantes que assinalaram as opções (D)Difícil e ( E ) Muito difícil.
Em cada barra foram assinalados também os extremos do intervalo de confiança de 95 \% como linhas verticais unidas por uma linha horizontal na forma da letra H maiúscula.
As Tabelas no Anexo II apresentam os valores absolutos e a distribuição percentual17 das alternativas válidas das nove questões, segundo o mesmo recorte de desempenho dos alunos e Grande Região de funcionamento do curso. 4.1 GRAU DE DIFICULDADE DA PROVA 4.1.1 Componente de Formação Geral Ao avaliarem Qual o grau de dificuldade desta prova na parte de Formação Geral? ( Questão 1 ), 23,2 \% do grupo de inscritos e presentes optaram pelas alternativas Difícil ou Muito difícil.
Entretanto, para mais da metade dos estudantes ( 60,9 \% ), o Componente de Formação Geral da prova foi considerado com grau de dificuldade Médio ( Gráfico 4.1, Gráfico 4.2 e, no Anexo II, a Tabela II.1 ).
O percentual de estudantes que consideraram a prova como Difícil ou Muito difícil foi maior na região Centro-Oeste, onde a proporção foi de 30,7 \%, enquanto a de menor incidência foi a Nordeste, com 16,8 \%.
No Gráfico 4.1, é possível observar que as diferenças entre a região Nordeste e as demais regiões são estatisticamente significativas, bem como entre a região Centro-Oeste e as demais.
Há ainda diferença estatisticamente significativa entre Sul e Sudeste.
Nas Grandes Regiões, a proporção de presentes à prova que consideraram o Componente de Formação Geral como sendo de grau de dificuldade Médio esteve entre 53,1 \% na região Centro-Oeste, e 66,8 \% na região Nordeste. 17 Cumpre lembrar uma das convenções para tabelas numéricas ( pág. iii ) sobre a possibilidade de a soma das partes não resultar em 100 \% por questões de arredondamento.
O percentual de alunos que consideraram a prova Difícil ou Muito difícil foi decrescente, em função dos quartos de desempenho : 33,9 \% no primeiro quarto e 13,0 \% no quarto, grupo de melhor desempenho na prova.
Nos quartos de desempenho intermediários, a proporção de alunos que consideraram a prova Difícil ou Muito difícil foi de 25,8 \% no segundo quarto e de 20,3 \% no terceiro quarto.
As diferenças entre os quartos são estatisticamente significativas.
Para todos os quartos de desempenho, a alternativa modal para essa pergunta foi Médio, com 55,5 \% e 62,8 \% dos respondentes nos quartos extremos, primeiro e quarto, respectivamente.
Gráfico 4.2 – Percentual de estudantes que avaliaram “ o grau de dificuldade desta prova na parte de Formação Geral ” como Difícil ou Muito difícil, por Quartos de Desempenho – Enade/2017 – Engenharia Mecânica Fonte : MEC/Inep/Daes - Enade/2017 4.1.2 Componente de Conhecimento Específico Ao responderem à Questão 2 – Qual o grau de dificuldade desta prova na parte de Componente Específico? – 46,6 \% do grupo de estudantes classificaram -na como Difícil ou Muito difícil.
Além disso, o Componente de Conhecimento Específico da prova foi considerado com grau de dificuldade Médio por 48,0 \% dos alunos ( Gráfico 4.3, Gráfico 4.4, e, no Anexo II, a Tabela II.2 ).
A análise das respostas dos estudantes quanto ao grau de dificuldade do Componente de Conhecimento Específico da prova, agregado por Grande Região, mostra que a diferença entre a menor e as maiores proporções de alunos que a avaliaram como Difícil ou Muito difícil é estatisticamente significativa, sendo a menor observada na região Nordeste ( 36,7 \% ).
As demais proporções foram : Sudeste ( 47,3 \% ), Sul ( 49,1 \% ) e Norte e Centro-Oeste ( 50,1 \% ).
O percentual de alunos que classificaram o grau de dificuldade como Médio, no Componente de Conhecimento Específico, variou de 43,2 \% a 58,0 \%, para as regiões Centro-Oeste e Nordeste, respectivamente.
Considerando -se a avaliação da dificuldade das questões do Componente de Conhecimento Específico da prova, de acordo com o desempenho dos estudantes, observa -se que há diferença estatisticamente significativa dos resultados entre o último e os demais quartos de desempenho, bem como entre o terceiro e os dois quartos inferiores.
A proporção dos que classificaram a parte específica como Difícil ou Muito difícil variou de 32,0 \% ( último quarto ) a 53,3 \% ( primeiro quarto ).
Como já comentado, a alternativa modal para a Questão 2 foi o grau Médio.
Já desagregando por quarto de desempenho, a situação é um pouco diferente com 42,1 \% do quarto inferior e 58,9 \% do superior optando por essa resposta. 4.2 EXTENSÃO DA PROVA EM RELAÇÃO AO TEMPO TOTAL Indagados quanto à extensão da prova, em relação ao tempo total oferecido para a sua resolução ( Questão 3 ), os estudantes apontaram, com maior incidência, a alternativa que considerava a extensão adequada, para todas as agregações consideradas ( Gráfico 4.5, Gráfico 4.6, e, no Anexo II, a Tabela II.3 ).
O percentual de alunos que responderam ser a extensão da prova adequada foi de 55,2 \%.
Já 35,0 \% dos inscritos presentes consideraram que a prova foi longa ou muito longa, e 9,8 \% a avaliaram como curta ou muito curta.
Dentre as Grandes Regiões, a proporção daqueles que avaliaram a prova como longa ou muito longa em relação ao tempo total destinado à sua resolução variou de 25,6 \% na região Norte até 39,4 \% na região Nordeste.
A diferença entre a região Norte e as demais, à exceção da Centro-Oeste, é estatisticamente significativa.
Também é estatisticamente significativa a diferença da região Nordeste e a das demais, à exceção da região Sul.
Considerando -se o desempenho dos alunos, nota -se que a mesma proporção ( 56,1 \% ) consideraram a extensão da prova adequada tanto no quarto de desempenho inferior, quanto no de melhor desempenho ( quarto superior ).
Nos quartos intermediários, essa proporção foi de 55,3 \% no segundo quarto e de 53,4 \% no terceiro.
No Gráfico 4.6, pode -se constatar que não há diferença estatisticamente significativa entre as proporções de escolha de estudantes que consideraram a prova longa ou muito longa, em função dos quartos de desempenho.
As proporções nos quartos extremos foram de 33,8 \% no primeiro quarto ( inferior ) e de 35,6 \% no quarto de melhor desempenho.
Nos quartos intermediários, essa proporção foi de 34,5 \% no segundo quarto e de 35,9 \% no terceiro. 4.3 COMPREENSÃO DOS ENUNCIADOS DAS QUESTÕES 4.3.1 Componente de Formação Geral Em relação aos enunciados das questões do Componente de Formação Geral ( Questão 4 ), as opiniões foram positivas, já que 77,4 \% dos alunos avaliados consideraram todos ou a maioria dos enunciados das questões claros e objetivos ( Gráfico 4.7, Gráfico 4.8, e, no Anexo II, a Tabela II.4 ).
Na análise regional, a percentagem de estudantes que avaliaram que todos ou a maioria dos enunciados das questões do Componente de Formação Geral estavam claros e objetivos variou de 75,8 \% na região Nordeste a 78,3 \% na região Sudeste.
Observa -se diferença estatisticamente significativa entre as regiões Sudeste e Sul.
A análise das percepções dos estudantes sobre a clareza e objetividade dos enunciados permite afirmar que todos, ou a maioria dos enunciados de questões relativas ao Componente de Formação Geral, foram considerados claros e objetivos para a maior parte dos respondentes ( maior ou igual a 75,8 \% em todas as regiões e maior ou igual a 72,6 \% para todos os quartos de desempenho ).
Segundo o desempenho, observa -se que a proporção dos que emitiram esta opinião variou entre 72,6 \% ( primeiro quarto ) e 80,3 \% ( terceiro quarto ), com diferenças estatisticamente significativas entre o quarto inferior e os demais quartos de desempenho, bem como entre os quartos intermediários.
No quarto superior, a clareza e objetividade de todos ou da maioria dos enunciados das questões foram percebidas por 79,6 \%. 4.3.2 Componente de Conhecimento Específico Em relação aos enunciados das questões do Componente de Conhecimento Específico da prova, para 81,6 \% dos estudantes avaliados da Área de Engenharia Mecânica, a clareza e a objetividade ( Questão 5 ) estavam presentes em todas ou na maioria das questões ( Gráfico 4.9, Gráfico 4.10, e no Anexo II, a Tabela II.5 ).
A maioria dos estudantes de todas as Grandes Regiões brasileiras considerou claros e objetivos todos ou a maioria dos enunciados das questões do Componente de Conhecimento Específico da prova, percentual sempre maior ou igual a 79,9 \%.
Não se observam diferenças estatisticamente significativas dentre as grandes Regiões.
A proporção de estudantes que consideraram os enunciados das questões claros e objetivos apresenta uma tendência crescente em relação ao aumento de desempenho : mais elevada no quarto superior ( 88,4 \% ) se comparada ao quarto inferior de desempenho ( 73,6 \% ).
As diferenças dentre todos os quartos de desempenho são estatisticamente significativas. 4.4 SUFICIÊNCIA DAS INFORMAÇÕES/INSTRUÇÕES FORNECIDAS Ao avaliarem as informações/instruções fornecidas para a resolução das questões ( Questão 6 ), 82,7 \% dos respondentes da Área de Engenharia Mecânica de todo o Brasil afirmaram que eram até excessivas ou suficientes em todas ou na maioria das questões ( Gráfico 4.11, Gráfico 4.12, e, no Anexo II, a Tabela II.6 ).
Quanto à distribuição de respondentes pelas Grandes Regiões observa -se que a proporção de estudantes que consideraram as informações/instruções fornecidas até excessivas ou suficientes em todas ou na maioria das questões foi sempre superior ou igual a 81,0 \%, chegando a 84,6 \% na região Nordeste.
Há diferença estatisticamente significativa da proporção observada na região Sudeste ( 81,9 \% ) em relação à das regiões Nordeste ( 84,6 \% ) e Sul ( 84,1 \% ).
Levando -se em conta o desempenho dos participantes, notam -se diferenças estatisticamente significativas entre as opiniões de estudantes de todos os quartos de desempenho, como mostra o Gráfico 4.12.
O percentual de participantes que avaliaram as informações/instruções como até excessivas ou suficientes em todas ou na maioria das questões foi mais elevado no quarto superior ( 91,2 \% ), percentual superior à média nacional ( 82,7 \% ).
Já no quarto inferior, a suficiência das informações/instruções declarada como até excessiva, em todas ou na maioria das questões foi percebida por 74,9 \% dos respondentes.
Essas proporções são crescentes com os quartos de desempenho. 4.5 DIFICULDADE ENCONTRADA AO RESPONDER À PROVA Indagados sobre as dificuldades com as quais se depararam ao responder à prova ( Questão 7 ), 16,2 \% dos estudantes apontaram o Desconhecimento do conteúdo.
Para 40,6 \%, a Forma diferente de abordagem do conteúdo foi indicada como dificuldade.
Já a Falta de motivação para fazer a prova foi a dificuldade apontada por 19,3 \% dos respondentes.
Considerando -se todo o Brasil, 14,0 \% dos respondentes afirmaram que não tiveram qualquer tipo de dificuldade para responder à prova ( Tabela II.7 no Anexo II ).
Os Gráficos 4.13 e 4.14 apresentam os percentuais de estudantes que apontaram o Desconhecimento do conteúdo como dificuldade ao responder à prova.
Na análise por Grandes Regiões, o percentual de inscritos e presentes que apontaram o Desconhecimento do conteúdo como dificuldade ao responder à prova não superou 18,6 \%.
Os percentuais variaram de 15,5 \% na região Sul a 18,6 \% na região Norte, não sendo observadas diferenças significativas estatisticamente dentre as regiões.
A Forma diferente de abordagem do conteúdo foi a escolha modal dos estudantes, com percentuais que variaram de 35,5 \% ( região Nordeste ) a 45,5 \% ( região Norte ).
O percentual de alunos que citaram a Falta de motivação como dificuldade variou de 9,4 \% ( região Norte ) a 23,4 \% ( região Nordeste ).
Os que declararam não ter qualquer tipo de dificuldade para responder à prova variaram de 12,3 \% na região Nordeste a 14,6 \% na região Sudeste.
Em relação aos quartos de desempenho, o Desconhecimento do conteúdo foi a opção escolhida por 17,5 \% dos estudantes do quarto superior e por 15,3 \% do quarto inferior.
Não se observa diferença estatisticamente significativa entre os quartos de desempenho.
A alternativa modal para os alunos, quando agregados pelos quartos de desempenho, foi que a dificuldade encontrada foi causada pela Forma diferente de abordagem do conteúdo : 46,1 \% no quarto inferior e 30,2 \% do quarto superior escolheram essa alternativa. 4.6 CONTEÚDOS DAS QUESTÕES OBJETIVAS DA PROVA Ao analisar os conteúdos das questões objetivas da prova ( Questão 8 ), um percentual muito pequeno dos estudantes avaliados, apenas 2,0 \%, afirmou que não estudou ainda a maioria desses conteúdos ( Gráficos 4.15, Gráfico 4.16, e a Tabelas II.8 no Anexo II ).
A maioria ( 72,2 \% ) afirmou ter estudado e aprendido muitos ou todos os conteúdos avaliados.
Na análise por Grande Região, a proporção de respondentes que escolheram a opção não estudou ainda a maioria desses conteúdos foi pequena.
Observa -se que na região Nordeste ( 3,1 \% ), apesar de pequena, a proporção foi maior do que a média nacional ( 2,0 \% ).
Observa -se diferença estatisticamente significativa entre a região Nordeste e as regiões Norte, Sudeste e Sul.
Em todas as regiões, a maioria dos presentes afirmou ter estudado e aprendido muitos ou todos os conteúdos, com proporções variando entre 65,3 \% na região Norte e 74,5 \% na Sul.
Considerando -se separadamente as opiniões de estudantes dos quatro quartos de desempenho, observa -se que, no quarto inferior, 4,2 \% ofereceram como resposta que não estudou ainda a maioria desses conteúdos, sendo 0,5 \% os do quarto superior com a mesma resposta.
A diferença entre os alunos que optaram por esse motivo de dificuldade entre todos os quartos de desempenho é estatisticamente significativa.
Levando -se em conta o quarto superior, 85,8 \% dos alunos afirmaram ter estudado e aprendido muitos ou todos os conteúdos.
No outro extremo, no primeiro quarto, 57,8 \% optaram pelas mesmas categorias.
As proporções são crescentes com o desempenho. 4.7 TEMPO GASTO PARA CONCLUIR A PROVA Ao responderem sobre o tempo de conclusão da prova ( Questão 9 ), quase três quartos dos estudantes ( 75,4 \% ) afirmaram ter gasto Entre duas e quatro horas ( Gráfico 4.17, Gráfico 4.18 e, no Anexo II, a Tabela II.9 ).
Considerando -se as cinco Grandes Regiões brasileiras, os percentuais dos que utilizaram Entre duas e quatro horas para finalizar a prova nas regiões Norte ( 74,7 \% ) e Nordeste ( 71,0 \% ) foram inferiores ao percentual nacional.
Nas demais Grandes Regiões, o percentual de alunos que dispensaram Entre duas e quatro horas para concluir a prova ficou igual ou acima de 75,7 \%, como mostrado no Gráfico 4.17.
Há diferenças estatisticamente significativas entre a proporção da região Nordeste e a das demais regiões, com exceção da região Norte.
Uma vez analisadas as alternativas escolhidas pelos estudantes que se situam nos diferentes quartos de desempenho, observa -se uma tendência crescente : uma proporção maior de participantes no quarto superior ( 80,8 \% ) declarou ter gasto Entre duas e quatro horas para concluir a prova, quando comparada com a dos quartos inferiores.
As diferenças entre o primeiro ( 69,8 \% ) e os demais quartos de desempenho são estatisticamente significativas, bem como as diferenças entre o quarto superior e os demais.

% matched lemmas: 
\end{textsample}
